%==============================================================================
% PACOTES ABNT — Livro: Odontologia & Inteligência Artificial
% Conforme NBR 14724 e NBR 6029
% Compilação: xelatex → bibtex → xelatex × 2
%==============================================================================

% --- Codificação e língua ---
\usepackage[utf8]{inputenc}
\usepackage{fontspec}
% babel já carregado pela classe abntex2
\usepackage{csquotes}
\usepackage{hyphenat}

% --- Geografia da página ABNT (NBR 14724) ---
\usepackage[a4paper,
  left=3cm,
  top=3cm,
  right=2cm,
  bottom=2cm
]{geometry}

% --- Matemática e ciência ---
\usepackage{amsmath, amssymb, amsthm}
\usepackage{siunitx}
\usepackage{textcomp}
\usepackage{gensymb}

% --- Cores ---
\usepackage{xcolor}
\definecolor{corPrimaria}{HTML}{1A5276}
\definecolor{corSecundaria}{HTML}{2E86C1}
\definecolor{corDestaque}{HTML}{E67E22}
\definecolor{corFundoCodigo}{HTML}{F5F5F5}
\definecolor{corBordaCodigo}{HTML}{CCCCCC}
\definecolor{corTesteSucesso}{HTML}{27AE60}
\definecolor{corTesteFalha}{HTML}{C0392B}
\definecolor{corNivelFacil}{HTML}{2ECC71}
\definecolor{corNivelMedio}{HTML}{F39C12}
\definecolor{corNivelDificil}{HTML}{E67E22}
\definecolor{corNivelAvancado}{HTML}{E74C3C}
\definecolor{corNivelPhD}{HTML}{8E44AD}
\definecolor{azulMedio}{HTML}{2E86C1}
\definecolor{verdeComentario}{HTML}{27AE60}
\definecolor{azPrimario}{HTML}{1A5276}
\definecolor{blueaccent}{HTML}{2E86C1}
\definecolor{codegreen}{HTML}{27AE60}
\definecolor{codegray}{HTML}{808080}
\definecolor{codeback}{HTML}{F5F5F5}
\definecolor{codeblue}{HTML}{2E86C1}
\definecolor{orangeaccent}{HTML}{E67E22}
\definecolor{cyanaccent}{HTML}{17A2B8}

% --- Figuras ---
\usepackage{graphicx}
\graphicspath{%
  {imagens/}%
  {imagens/figuras/cap01/}{imagens/figuras/cap02/}%
  {imagens/figuras/cap03/}{imagens/figuras/cap04/}%
  {imagens/figuras/cap05/}{imagens/figuras/cap06/}%
  {imagens/figuras/cap07/}{imagens/figuras/cap08/}%
  {imagens/figuras/cap09/}{imagens/figuras/cap10/}%
  {imagens/figuras/cap11/}{imagens/figuras/cap12/}%
  {imagens/figuras/cap13/}{imagens/figuras/cap14/}%
  {imagens/figuras/cap15/}%
}
\usepackage[export]{adjustbox}
\usepackage{caption}
\usepackage{subcaption}
\usepackage{float}

% --- TikZ e gráficos vetoriais ---
\usepackage{tikz}
\usetikzlibrary{shapes, arrows, positioning, calc, decorations.pathreplacing}
\usepackage{pgfplots}
\pgfplotsset{compat=1.18}

% --- Tabelas ---
\usepackage{booktabs}
\usepackage{longtable}
\usepackage{tabularx}
\usepackage{array}
\usepackage{multirow}
\usepackage{colortbl}
\usepackage{ltablex}

% --- Código-fonte (listings) ---
\usepackage{listings}
\lstset{
  language=Python,
  basicstyle=\footnotesize\ttfamily,
  keywordstyle=\bfseries\color{azulMedio},
  commentstyle=\itshape\color{verdeComentario},
  stringstyle=\color{magenta},
  numberstyle=\tiny\color{gray},
  numbers=left, numbersep=5pt,
  tabsize=2, breaklines=true,
  prebreak=\raisebox{0ex}[0ex][0ex]{\ensuremath{\hookleftarrow}},
  frame=single, frameround=tttt,
  rulecolor=\color{corBordaCodigo},
  backgroundcolor=\color{corFundoCodigo},
  showstringspaces=false, showtabs=false,
  xleftmargin=2em, framexleftmargin=2em,
  literate={°}{{\degree}}1,
  morekeywords={True, False, None, self, class, def, return, yield, import,
    from, as, with, assert, async, await, print, len, range, enumerate, zip,
    map, filter, lambda, in, not, and, or, is, if, elif, else, for, while, try,
    except, finally, raise, break, continue, pass},
  emphstyle=\color{azPrimario}
}
\lstnewenvironment{saida}[1][]{%
  \lstset{language={},basicstyle=\footnotesize\ttfamily\color{black},
    backgroundcolor=\color{black!5},frame=single,numbers=none,#1}}{}

% Estilo 'saida' (para \begin{lstlisting}[style=saida]) — espelha o ambiente saida
\lstdefinestyle{saida}{language={},basicstyle=\footnotesize\ttfamily\color{black},
  backgroundcolor=\color{black!5},frame=single,numbers=none}

% Linguagens ausentes no listings (usadas nos capítulos de engenharia)
\lstdefinelanguage{yaml}{%
  keywords={true,false,null,yes,no,on,off},
  sensitive=true, comment=[l]{\#},
  morestring=[b]', morestring=[b]"}
\lstdefinelanguage{docker}{%
  keywords={FROM,RUN,CMD,LABEL,MAINTAINER,EXPOSE,ENV,ADD,COPY,ENTRYPOINT,%
    VOLUME,USER,WORKDIR,ARG,ONBUILD,STOPSIGNAL,HEALTHCHECK,SHELL},
  sensitive=true, comment=[l]{\#}, morestring=[b]"}

% --- Caixas destacadas ---
\usepackage{mdframed}

% --- Camada imersiva (modelo sensorial) ---
\usepackage[most]{tcolorbox}
% sidenotes exige changepage[strict]; abntex2/memoir já "emula" changepage
% sem opções e define \sidecaption/marginfigure/margintable — resolvemos o
% choque de opções/nomes ANTES de carregar sidenotes (nenhum desses comandos
% do memoir é usado nos capítulos do livro; sidenotes assume o papel deles).
\PassOptionsToPackage{strict}{changepage}
\makeatletter
\let\sidecaption\undefined
\let\marginfigure\undefined
\let\endmarginfigure\undefined
\let\margintable\undefined
\let\endmargintable\undefined
\makeatother
\usepackage{sidenotes}   % requer/puxa marginnote
\usepackage{fontawesome5}
\usepackage{qrcode}

% --- Aumentar profundidade máxima de aninhamento de listas ---
\usepackage{enumitem}
\setlist{maxdepth=6}
% Também aumentar os contadores padrão do LaTeX
\renewcommand{\theenumi}{\arabic{enumi}}
\renewcommand{\theenumii}{\alph{enumii}}
\renewcommand{\theenumiii}{\roman{enumiii}}
\renewcommand{\theenumiv}{\Alph{enumiv}}

% --- Hiperlinks ---
% hyperref já carregado pela classe abntex2
\hypersetup{
  bookmarksnumbered=true, bookmarksopen=true, bookmarksopenlevel=1,
  colorlinks=true, linkcolor=black, citecolor=black, urlcolor=black,
  linktoc=all,
  pdfauthor={\LivroAutor},
  pdftitle={\LivroTitulo\ -- \LivroSubtitulo},
  pdfsubject={Odontologia e Inteligência Artificial},
  pdfkeywords={\LivroPalavrasChave}
}
\usepackage{bookmark}

% --- Índices e glossário ---
\usepackage[toc, acronym, nomain, style=alttree]{glossaries}
\makeglossaries
\usepackage{makeidx}
\makeindex

% --- Referências ABNT (NBR 6023) ---
\usepackage[backend=biber, style=abnt-numeric, giveninits=true, uniquename=init]{biblatex}
\addbibresource{referencias/bibliografia.bib}
\addbibresource{referencias/dois-pendentes.bib}

% --- Miscelânea ---
\usepackage[bottom]{footmisc}
\usepackage{pdfpages}
\usepackage{tocloft}
\usepackage{etoolbox}
\usepackage{ifplatform}
\usepackage{datetime2}
\usepackage[type={CC}, modifier={by-nc-sa}, version={4.0}]{doclicense}
\usepackage{lettrine}   % letra capitular (drop cap) por nível

% --- Listas de figuras e tabelas (memoir/abntex2) ---
\renewcommand{\cftfigurefont}{Figura }
\renewcommand{\cfttablefont}{Tabela }
\addto\captionsbrazil{\renewcommand{\listfigurename}{Lista de Figuras}}
\addto\captionsbrazil{\renewcommand{\listtablename}{Lista de Tabelas}}
